\chapter{The Encoder and the Verification Lift}\label{ch:encoder}

\NewV\ \emph{Two roles the architecture cannot do without and has not yet
derived. This chapter exists to keep them visible as open problems rather
than letting them be quietly filled by an import.}

\begin{workplan}{\NewV}
\wpgoal{State precisely what the encoder and the verification lift must
\emph{do} --- as constraints the rest of the mathematics imposes on them ---
without assuming any particular architecture for either.}

\wpsource{\texttt{TATIANA\_V1\_IDENTITY\_MAP.md}~\S2.1 (the narrow allowance,
now further narrowed --- see \S\ref{sec:encoder-position} of
Chapter~\ref{ch:v2}); \texttt{MATH\_TOOLBOX.md}~I.8 (the verification
trichotomy, already fixed epistemically), I.5 (what the drive term needs to
be for the wake flow to act on it).}

\wpsteps{
  \item \textbf{Derive the encoder's type, do not choose it.} Work backwards
    from the wake phase: the drive term enters $\dot s = -\Lap_{\Fs}s +
    \text{drive}$, so $o_t$ must live where the restriction maps can receive
    it. Establish the minimum: which space, what normalisation, and whether
    anything more than a vector is required. The answer is a \emph{constraint},
    and the chapter's contribution is to state it exactly.
  \item \textbf{Ask what would make an encoder wrong.} Two failure modes to
    develop: an encoder whose output geometry already contains the structure
    the sheaf is supposed to learn (in which case the sheaf is decorative and
    the architecture is a dressed-up embedding lookup), and an encoder whose
    output is isotropic noise to the sheaf (in which case nothing can ever
    glue). Both are checkable. Specify how.
  \item \textbf{Derive the verification lift's contract.} What must
    $\mathcal{L}$ guarantee for the trichotomy --- \textsc{verified},
    \textsc{refuted}, \textsc{unverifiable} --- to remain meaningful? The
    third value is the load-bearing one: it exists precisely so that
    inability to check is never recorded as falsehood, and a checker crash
    must return it rather than a refutation. State the contract as
    requirements on the lift, not as a description of a particular checker.
  \item \textbf{Connect the contract to consolidation.} $\gamma(\nu)$ reads
    the verdict: $\gamma_0$ if verified, $\varepsilon\gamma_0$ if
    unverifiable, $0$ if refuted. So the lift's guarantees are exactly what
    stops unverified structure from entering the invariant. Make the
    dependency explicit --- it is why this chapter cannot be deferred
    indefinitely.
  \item \textbf{Check the surviving toolbox before reaching outward.} Ask
    whether either role reduces to something already present: is an encoder a
    restriction map into the complex from a boundary cell? Is a checker a
    discrete verification map with the same signature as the operators the
    book already defines? Record the answers even if negative, since a
    negative answer is what would justify building something new.
  \item \textbf{State the standing constraint, in the chapter, not in a
    footnote.} Neither role may be discharged by wrapping a pretrained,
    externally-hosted model. If a derivation arrives at machinery that also
    appears in such models, that is permitted --- as a conclusion with a
    derivation attached, never as a default and never as an import.
}

\wpdone{Both roles have stated contracts derived from what the rest of the
architecture requires of them; at least two failure modes are specified as
checkable; and no architecture has been assumed for either.}
\end{workplan}
